\documentclass[11pt]{article}
\usepackage[T1]{fontenc}
\usepackage[utf8]{inputenc}
\usepackage{mathptmx}
\usepackage{microtype}
\usepackage[margin=2.5cm]{geometry}
\usepackage{hyperref}
\pagestyle{empty}

\begin{document}

\noindent
The Editors\\
\textit{B.E.\ Journal of Economic Analysis \& Policy}\\
De Gruyter

\bigskip
\noindent
London, 16 May 2026

\bigskip
\noindent
Dear Editors,

\medskip

I write to submit a revised version of the manuscript
``Participation Costs, Social Capital, and the Limits of Mobility
Policy: A Model of Status Competition under Income Inequality''
for consideration in \textit{BEJEAP}. This follows an earlier version 
submitted in October 2025 that was rejected. The current version
responds directly to the comments made by the referee and constitutes a significant revision.

The paper employs microeconomic theory to analyse when and why
income inequality suppresses participation in status competitions
for elite credentials, and what policy instruments can restore
mobility. The paper situates itself in the positional competition tradition
rather than the human capital tradition - viewing social capital as a
rank-building input into status tournaments rather than
informational. The main findings of the current version are a hard mobility ceiling
$\mu/[2(1-\mu)]$ determined by the design of the
competition rather than by who participates in it, and a formal
demonstration that participation-cost reduction and competition
redesign are complementary instruments. I believe that the results speak directly to the design of
credential programmes, hiring processes, and access policies.

The two major issues identified by the referee were: the paper's restriction to
$N=2$ players and the absence of clear propositions. The first issue has now been addressed by letting the model place one high-income participant against a general $Q \geq 1$ low-income challengers and by deriving all results --- winning
probabilities, Nash equilibrium conditions, the participation
threshold, and the mobility ceiling --- for a general $Q$. The second issue has been addressed by adding formal propositions with proofs for the Case~II equilibrium (Proposition~1) and the asymptotic probability surplus (Proposition~2). A new Section~3 now also develops the policy content at length.

The referee had also raised four other concerns -- i) interpretations from a zero-sum structure ii) a misleading use of ``signaling'' iii) an insufficient engagement with existing literature and iv) restriction to uniform distributions in the model. To address the zero-sum structural limitations, the paper now explicitly limits its scope to positional competitions with a fixed number of elite positions ---
partnership tracks, academic appointments, competitive admissions
--- where the zero-sum characterisation is accurate rather than
approximate. To correct the misuse of ``signaling'', all language around signaling has now been removed and the paper's mechanism is stated as contest-theoretic rather than signaling-theoretic. The mechanism used in the model is characterised explicitly as a productive rank-building input in a tournament - where social capital enters the composite score directly and raises the probability of winning through a contest scoring rule. The paper now engages with Gallice and Grillo (2019) on two aspects: contest versus signaling mechanism; and the threshold cost structure that generates binary
equilibrium outcomes for formulation of a mobility ceiling. To address the limitations from uniform distributions, a robustness remark has been added -- explaining that the key qualitative results i.e. the binary equilibrium structure and all comparative statics, hold for any pair of symmetric distributions on $[0,1]$. The
closed-form expressions are still specific to the uniform case - but only as a 
tractable parametric representation of a more general result.


I confirm that this manuscript is not under consideration
elsewhere, that the paper has not been published previously,
and that I declare no conflicts of interest.

\bigskip
\noindent
Yours sincerely,

\bigskip
\noindent
Anurag Srivastava\\

\end{document}
